\begin{figurebox}[Courbe de $f(x)=(x-3)^2-4$. Les deux racines et le sommet sont indiqués ; l’axe de symétrie est la droite $x=3$.]
\begin{tikzpicture}\begin{axis}[width=.88\linewidth,height=6.4cm,axis lines=middle,grid=major,grid style={gray!25},tick label style={font=\normalsize},label style={font=\normalsize},clip=false,xmin=-.5,xmax=6.5,ymin=-5,ymax=7,xlabel={$x$},ylabel={$f(x)$},xtick={0,1,2,3,4,5,6},ytick={-4,-2,0,2,4,6}]\addplot[mxink,very thick,domain=0:6,samples=100]{(x-3)^2-4};\draw[dashed,gray] (axis cs:3,-4)--(axis cs:3,6);\addplot[only marks,mark=*,mxmint] coordinates {(1,0)(5,0)(3,-4)};\node[above left,fill=white,inner sep=2pt]at(axis cs:1,0){$\pt{1;0}$};\node[above right,fill=white,inner sep=2pt]at(axis cs:5,0){$\pt{5;0}$};\node[below right,fill=white,inner sep=2pt]at(axis cs:3,-4){$S\pt{3;-4}$};\end{axis}\end{tikzpicture}
\end{figurebox}
